\documentclass[10pt,a4paper]{amsart}
\usepackage{amsmath,amssymb,amsthm,graphicx}
% Retained minimal preamble for the Stacks Project chapter more-algebra.tex (More on Algebra).
% The chapter file pinned-original-source.tex carries only \input{preamble} on its first line, so no
% preamble text lives inside the chapter itself.  The Stacks Project uses one single shared preamble for
% every chapter; the snapshot stages it identically as sources/stacks-chow/preamble.tex and
% sources/stacks-etale-cohomology/preamble.tex.  The block below is copied mechanically, line for line,
% from sources/stacks-chow/preamble.tex lines 14-18 (the Xy-pic setup) and 154-217 (the native theorem
% declarations together with the mathematical macros).  Nothing was added, removed or redefined; the
% wrapper class and the amsmath/amssymb/amsthm/graphicx base packages are supplied by the assembler.
% Chapter: more-algebra.tex -- More on Algebra.
% For commutative diagrams we use Xy-pic
\usepackage[all]{xy}

% We use 2cell for 2-commutative diagrams.
\xyoption{2cell}

% Theorem environments.
%
\theoremstyle{plain}
\newtheorem{theorem}[subsection]{Theorem}
\newtheorem{proposition}[subsection]{Proposition}
\newtheorem{lemma}[subsection]{Lemma}

\theoremstyle{definition}
\newtheorem{definition}[subsection]{Definition}
\newtheorem{example}[subsection]{Example}
\newtheorem{exercise}[subsection]{Exercise}
\newtheorem{situation}[subsection]{Situation}

\theoremstyle{remark}
\newtheorem{remark}[subsection]{Remark}
\newtheorem{remarks}[subsection]{Remarks}

\numberwithin{equation}{subsection}

% Macros
%
\def\lim{\mathop{\mathrm{lim}}\nolimits}
\def\colim{\mathop{\mathrm{colim}}\nolimits}
\def\Spec{\mathop{\mathrm{Spec}}}
\def\Hom{\mathop{\mathrm{Hom}}\nolimits}
\def\Ext{\mathop{\mathrm{Ext}}\nolimits}
\def\SheafHom{\mathop{\mathcal{H}\!\mathit{om}}\nolimits}
\def\SheafExt{\mathop{\mathcal{E}\!\mathit{xt}}\nolimits}
\def\Sch{\mathit{Sch}}
\def\Mor{\mathop{\mathrm{Mor}}\nolimits}
\def\Ob{\mathop{\mathrm{Ob}}\nolimits}
\def\Sh{\mathop{\mathit{Sh}}\nolimits}
\def\NL{\mathop{N\!L}\nolimits}
\def\CH{\mathop{\mathrm{CH}}\nolimits}
\def\proetale{{pro\text{-}\acute{e}tale}}
\def\etale{{\acute{e}tale}}
\def\QCoh{\mathit{QCoh}}
\def\Ker{\mathop{\mathrm{Ker}}}
\def\Im{\mathop{\mathrm{Im}}}
\def\Coker{\mathop{\mathrm{Coker}}}
\def\Coim{\mathop{\mathrm{Coim}}}

% Boxtimes
%
\DeclareMathSymbol{\boxtimes}{\mathbin}{AMSa}{"02}

%
% Macros for moduli stacks/spaces
%
\def\QCohstack{\mathcal{QC}\!\mathit{oh}}
\def\Cohstack{\mathcal{C}\!\mathit{oh}}
\def\Spacesstack{\mathcal{S}\!\mathit{paces}}
\def\Quotfunctor{\mathrm{Quot}}
\def\Hilbfunctor{\mathrm{Hilb}}
\def\Curvesstack{\mathcal{C}\!\mathit{urves}}
\def\Polarizedstack{\mathcal{P}\!\mathit{olarized}}
\def\Complexesstack{\mathcal{C}\!\mathit{omplexes}}
% \Pic is the operator that assigns to X its picard group, usage \Pic(X)
% \Picardstack_{X/B} denotes the Picard stack of X over B
% \Picardfunctor_{X/B} denotes the Picard functor of X over B
\def\Pic{\mathop{\mathrm{Pic}}\nolimits}
\def\Picardstack{\mathcal{P}\!\mathit{ic}}
\def\Picardfunctor{\mathrm{Pic}}
\def\Deformationcategory{\mathcal{D}\!\mathit{ef}}
\usepackage{cleveref}
\crefname{originallemma}{original lemma}{original lemmas}
\Crefname{originallemma}{Original lemma}{Original lemmas}
\makeatletter
\newlabel{more-algebra-lemma-formal-fibres-Rk}{{0BIZ}{source}{}{sourcepointer-1}{}}
\newlabel{more-algebra-lemma-formal-fibres-Rk@cref}{{[originallemma][0][]0BIZ}{[1][0][]source}{}{}{}}
\makeatother
\crefname{originallemma}{original lemma}{original lemmas}
\Crefname{originallemma}{Original lemma}{Original lemmas}
\makeatletter
\newlabel{algebra-lemma-make-separable}{{030R}{source}{}{sourcepointer-2}{}}
\newlabel{algebra-lemma-make-separable@cref}{{[originallemma][0][]030R}{[1][0][]source}{}{}{}}
\makeatother
\crefname{originallemma}{original lemma}{original lemmas}
\Crefname{originallemma}{Original lemma}{Original lemmas}
\makeatletter
\newlabel{algebra-lemma-localization-smooth-separable}{{037X}{source}{}{sourcepointer-3}{}}
\newlabel{algebra-lemma-localization-smooth-separable@cref}{{[originallemma][0][]037X}{[1][0][]source}{}{}{}}
\makeatother
\crefname{originallemma}{original lemma}{original lemmas}
\Crefname{originallemma}{Original lemma}{Original lemmas}
\makeatletter
\newlabel{algebra-lemma-base-change-smooth}{{00T4}{source}{}{sourcepointer-4}{}}
\newlabel{algebra-lemma-base-change-smooth@cref}{{[originallemma][0][]00T4}{[1][0][]source}{}{}{}}
\makeatother
\crefname{originallemma}{original lemma}{original lemmas}
\Crefname{originallemma}{Original lemma}{Original lemmas}
\makeatletter
\newlabel{algebra-lemma-Rk-goes-up}{{033A}{source}{}{sourcepointer-5}{}}
\newlabel{algebra-lemma-Rk-goes-up@cref}{{[originallemma][0][]033A}{[1][0][]source}{}{}{}}
\makeatother
\crefname{originallemma}{original lemma}{original lemmas}
\Crefname{originallemma}{Original lemma}{Original lemmas}
\makeatletter
\newlabel{algebra-lemma-characterize-smooth-over-field}{{00TT}{source}{}{sourcepointer-6}{}}
\newlabel{algebra-lemma-characterize-smooth-over-field@cref}{{[originallemma][0][]00TT}{[1][0][]source}{}{}{}}
\makeatother
\crefname{originallemma}{original lemma}{original lemmas}
\Crefname{originallemma}{Original lemma}{Original lemmas}
\makeatletter
\newlabel{algebra-lemma-descent-Rk}{{0353}{source}{}{sourcepointer-7}{}}
\newlabel{algebra-lemma-descent-Rk@cref}{{[originallemma][0][]0353}{[1][0][]source}{}{}{}}
\makeatother
\crefname{originallemma}{original lemma}{original lemmas}
\Crefname{originallemma}{Original lemma}{Original lemmas}
\makeatletter
\newlabel{algebra-lemma-separable-algebraic-diagonal}{{0C2X}{source}{}{sourcepointer-8}{}}
\newlabel{algebra-lemma-separable-algebraic-diagonal@cref}{{[originallemma][0][]0C2X}{[1][0][]source}{}{}{}}
\makeatother
\title{Regularity of formal fibres for a smooth algebra over a field}
\author{The Stacks Project Authors}
\date{stacks-project:more-algebra.tex:local-snapshot-20261012; GNU FDL 1.2 or later}
\begin{document}
\maketitle
\setcounter{subsection}{1}
\noindent\textit{Editorial note.} All delivered mathematical text is the author's own, copied from the pinned Stacks Project chapter file more-algebra.tex; no statement or proof step was rewritten, reordered or removed.\par\medskip
\noindent\textit{Editorial note.} The chapter file pinned-original-source.tex carries only \textbackslash{}input\{preamble\} on its first line, so no preamble text lives inside the chapter. The retained file sources/stacks-more-algebra/retained-preamble.tex is a mechanical, line-for-line copy of the project's single shared preamble as staged in the snapshot (Xy-pic setup plus theorem declarations and mathematical macros). Nothing was added, removed or redefined, and the excerpt serves only as preamble, never as a span coordinate.\par\medskip
\noindent\textit{Editorial note.} Both counted spans start at the \textbackslash{}begin line of their environment, so the statement and the author's proof are delivered whole; surrounding prose of the section is deliberately outside the counted spans and is carried as standard prerequisites instead.\par\medskip
\noindent\textit{Editorial note.} The citation fields (source\_id, source\_version, authors, source\_url, license.name, source\_title) are identical to those of every other card this lane has produced from this chapter, so the batch selector treats them as one source.\par\medskip
\noindent\textit{Allowed standard prerequisites (editorial).}\begin{itemize}
\item Regularity of a local ring and its preservation under etale maps and localisation. used at proof lines 13276-13320. Carries the property from the polynomial case to the smooth case.
\item Structure of a smooth algebra over a field as a localisation of an etale algebra over a polynomial ring. used at proof lines 13300-13320. Reduces the statement to the polynomial case.
\item The completion of a polynomial ring at a maximal ideal is a power series ring with as many variables as generators. used at proof lines 13320-13332. Is the explicit computation that supplies regularity in the remaining case.
\end{itemize}\medskip
\ifdefined\hypertarget\hypertarget{sourcepointer-1}{}\fi
\noindent\textit{Source reference.} Original lemma 0BIZ: focused-seeds/AGP/sources/stacks-snapshot-20261011/tags.txt, source lines 14507--14507. The lemma's own canonical tag, quoted so that readers can resolve the label inside the delivered text; it points at the counted statement itself and carries no proof content.\par\medskip
\ifdefined\hypertarget\hypertarget{sourcepointer-2}{}\fi
\noindent\textit{Source reference.} Original lemma 030R: focused-seeds/AGP/sources/stacks-snapshot-20261011/tags.txt, source lines 3958--3958. Navigation pointer to a statement the counted proof refers to but does not contain; nothing load-bearing is delegated to it.\par\medskip
\ifdefined\hypertarget\hypertarget{sourcepointer-3}{}\fi
\noindent\textit{Source reference.} Original lemma 037X: focused-seeds/AGP/sources/stacks-snapshot-20261011/tags.txt, source lines 4216--4216. Navigation pointer to a statement the counted proof refers to but does not contain; nothing load-bearing is delegated to it.\par\medskip
\ifdefined\hypertarget\hypertarget{sourcepointer-4}{}\fi
\noindent\textit{Source reference.} Original lemma 00T4: focused-seeds/AGP/sources/stacks-snapshot-20261011/tags.txt, source lines 1062--1062. Navigation pointer to a statement the counted proof refers to but does not contain; nothing load-bearing is delegated to it.\par\medskip
\ifdefined\hypertarget\hypertarget{sourcepointer-5}{}\fi
\noindent\textit{Source reference.} Original lemma 033A: focused-seeds/AGP/sources/stacks-snapshot-20261011/tags.txt, source lines 4049--4049. Navigation pointer to a statement the counted proof refers to but does not contain; nothing load-bearing is delegated to it.\par\medskip
\ifdefined\hypertarget\hypertarget{sourcepointer-6}{}\fi
\noindent\textit{Source reference.} Original lemma 00TT: focused-seeds/AGP/sources/stacks-snapshot-20261011/tags.txt, source lines 1087--1087. Navigation pointer to a statement the counted proof refers to but does not contain; nothing load-bearing is delegated to it.\par\medskip
\ifdefined\hypertarget\hypertarget{sourcepointer-7}{}\fi
\noindent\textit{Source reference.} Original lemma 0353: focused-seeds/AGP/sources/stacks-snapshot-20261011/tags.txt, source lines 4114--4114. Navigation pointer to a statement the counted proof refers to but does not contain; nothing load-bearing is delegated to it.\par\medskip
\ifdefined\hypertarget\hypertarget{sourcepointer-8}{}\fi
\noindent\textit{Source reference.} Original lemma 0C2X: focused-seeds/AGP/sources/stacks-snapshot-20261011/tags.txt, source lines 15172--15172. Navigation pointer to a statement the counted proof refers to but does not contain; nothing load-bearing is delegated to it.\par\medskip
\begin{lemma}
\label{lemma-formal-fibres-Rk}
Fix $n \geq 0$. Properties (A), (B), (C), (D), and (E) hold for
$P(k \to R) =$``$R \otimes_k k'$ has $(R_n)$ for all finite
extensions $k'/k$''.
\end{lemma}

\begin{proof}
Let $k \to R$ be a ring map where $k$ is a field and $R$ a Noetherian
ring. Assume $P(k \to R)$ is true.
Let $K/k$ be a finitely generated field extension.
By Algebra, Lemma \ref{algebra-lemma-make-separable} we can find a diagram
$$
\xymatrix{
K \ar[r] & K' \\
k \ar[u] \ar[r] & k' \ar[u]
}
$$
where $k'/k$, $K'/K$ are finite purely inseparable field
extensions such that $K'/k'$ is separable. By
Algebra, Lemma \ref{algebra-lemma-localization-smooth-separable}
there exists a smooth $k'$-algebra $B$ such that $K'$ is the
fraction field of $B$. Now we can argue as follows:
Step 1: $R \otimes_k k'$ satisfies $(R_n)$ because we assumed $P$
for $k \to R$.
Step 2: $R \otimes_k k' \to R \otimes_k k' \otimes_{k'} B$ is a
smooth ring map (Algebra, Lemma \ref{algebra-lemma-base-change-smooth})
and we conclude $R \otimes_k k' \otimes_{k'} B$ satisfies $(R_n)$
by Algebra, Lemma \ref{algebra-lemma-Rk-goes-up}
(and using Algebra, Lemma \ref{algebra-lemma-characterize-smooth-over-field}
to see that the hypotheses are satisfied).
Step 3. $R \otimes_k k' \otimes_{k'} K' = R \otimes_k K'$ satisfies
$(R_n)$ as it is a localization of a ring having $(R_n)$.
Step 4. Finally $R \otimes_k K$ satisfies $(R_n)$ by descent of
$(R_n)$ along the faithfully flat ring map
$K \otimes_k R \to K' \otimes_k R$
(Algebra, Lemma \ref{algebra-lemma-descent-Rk}).
This proves (A).
Part (B) follows too: a Noetherian ring has $(R_n)$ if and only if
all of its local rings have $(R_n)$.
Part (C). This follows from Algebra, Lemma \ref{algebra-lemma-Rk-goes-up}
as the fibres of a regular homomorphism are regular (small detail omitted).
Part (D). This follows from Algebra, Lemma \ref{algebra-lemma-descent-Rk}
(small detail omitted).

\medskip\noindent
Part (E). Let $l/k$ be a separable algebraic extension of fields and
let $l \to R$ be a ring map with $R$ Noetherian. Assume that
$k \to R$ has $P$. We have to show that $l \to R$ has $P$.
Let $l'/l$ be a finite extension. First observe that there exists a
finite subextension $l/m/k$ and a finite extension $m'/m$
such that $l' = l \otimes_m m'$. Then $R \otimes_l l' = R \otimes_m m'$.
Hence it suffices to prove that $m \to R$ has property $P$, i.e.,
we may assume that $l/k$ is finite. If $l/k$ is finite, then $l'/k$
is finite and we see that
$$
l' \otimes_l R = (l' \otimes_k R) \otimes_{l \otimes_k l} l
$$
is a localization
(by Algebra, Lemma \ref{algebra-lemma-separable-algebraic-diagonal})
of the Noetherian ring $l' \otimes_k R$ which has property $(R_n)$
by assumption $P$ for $k \to R$. This proves that $l' \otimes_l R$
has property $(R_n)$ as desired.
\end{proof}

\end{document}
